\phantomsection
\label{fig:han-period-timeline}
\begin{center}
\begin{tikzpicture}[x=.035cm,y=1cm,font=\small]
  \fill[paper] (-218,-.6) rectangle (232,2.55);
  \draw[very thick,dynasty] (-206,1.25) -- (220,1.25);

  \foreach \year/\labeltext/\side in {
    -206/{前206\\分封}/above,
    -202/{前202\\汉立}/below,
    -154/{前154\\七国}/above,
    -141/{前141\\武帝}/below,
    -119/{前119\\漠北}/above,
    -91/{前91\\太子危机}/below,
    -9/{前9\\王莽居摄}/above,
    9/{9\\新}/below,
    23/{23\\更始}/above,
    25/{25\\复汉}/below,
    89/{89\\北边纪功}/above,
    184/{184\\黄巾}/below,
    189/{189\\宫门失守}/above,
    200/{200\\官渡}/below,
    208/{208\\赤壁}/above,
    220/{220\\禅代}/below}
  {
    \draw[timeline,thick] (\year,1.12) -- (\year,1.38);
    \node[\side,align=center,font=\scriptsize,text=ink] at (\year,1.42) {\labeltext};
  }

  \draw[dashed,source!75] (9,.08) -- (9,2.3);
  \draw[dashed,source!75] (25,.08) -- (25,2.3);
  \node[align=center,font=\scriptsize,text=dynasty] at (-100,.15) {西汉：关中、郡国、边疆与财政的重组};
  \node[align=center,font=\scriptsize,text=timeline] at (17,2.27) {新与复汉：命令、名分和执行的断裂};
  \node[align=center,font=\scriptsize,text=dynasty] at (124,.15) {东汉：雒阳复接、地方军权与区域竞争};

  \node[draw=source!75,fill=paper,rounded corners=2pt,align=left,
    inner sep=3pt,font=\scriptsize] at (105,2.45)
    {时间轴为编年阅读的定位线，不表示事件之间\\
     存在自动因果，也不按战争规模或政权控制范围绘制。};
\end{tikzpicture}

\smallskip
\small\textit{图\quad 汉帝国时期时间轴（前206--220）：据《史记》《汉书》《后汉书》《后汉纪》《三国志》与《资治通鉴》的可核年序编制。它只标出本卷反复回看的转换点；政权更替、战争与制度动作的精确日期、参与者和影响范围仍以各节的来源注记为准。\quad\hyperref[sec:han-shared-chronology]{回到共享年序}}
\end{center}
